% =============================================================================
% GLOSSÁRIO HISTÓRICO E PERICIAL — APÊNDICE DE APOIO AO LEITOR
% =============================================================================
\cleardoublepage
\thispagestyle{empty}
\chapter*{Glossário Histórico e Pericial}
\addcontentsline{toc}{chapter}{Glossário Histórico e Pericial}
\markboth{Glossário Histórico e Pericial}{}

\begingroup

\noindent\textit{Vocabulário de apoio para navegação nos documentos da Seca de 1915 e dos laudos psiquiátricos do Hospital Colônia de Barbacena (1917--1979).}

\vspace{0.8em}
\noindent\rule{\textwidth}{0.8pt}
\vspace{0.6em}

\noindent\textbf{1. TERMOS DA SECA E DO SERTÃO (1914--1915)}

\vspace{0.4em}
\begin{itemize}[itemsep=0.2em]
  \item \textbf{Campo do Alagadiço (1915):} Campo de confinamento de retirantes criado próximo a Fortaleza durante a Seca de 1915. É o referente histórico adotado para o cenário ficcional da primeira parte deste livro.
  \item \textbf{Campo do Patu (1932--1933):} Campo situado em Senador Pompeu e associado à seca de 1932. Não é o campo de 1915 narrado no percurso de Joaquim.
  \item \textbf{Retirantes / Molambudos:} Designação popular pejorativa dada aos sertanejos que fugiam do interior a pé, em farrapos de pano (molambos), tentando sobreviver à estiagem extrema.
  \item \textbf{Vala Comum:} Cova coletiva rasa onde centenas de corpos de vítimas da fome e do tifo eram depositados diariamente sem identificação nos campos de concentração de 1915.
\end{itemize}

\vspace{0.8em}
\noindent\rule{\textwidth}{0.5pt}
\vspace{0.6em}

\noindent\textbf{2. TERMOS MÉDICOS E PSIQUIÁTRICOS DO SÉCULO XX}

\vspace{0.4em}
\begin{itemize}[itemsep=0.2em]
  \item \textbf{Hospital Colônia de Barbacena:} Instituição psiquiátrica fundada em Minas Gerais em 1903. Entre 1930 e 1980, tornou-se centro de internamento em massa de marginalizados sem diagnóstico psiquiátrico formal, resultando na morte de mais de 60.000 pessoas (episódio conhecido como \textit{Holocausto Brasileiro}).
  \item \textbf{Eletroconvulsoterapia (ECT / Eletrochoque):} Aplicação de corrente elétrica no cérebro para induzir convulsões. No Hospital Colônia, era praticada de forma contínua e punitiva, sem anestesia ou gel condutor, provocando queimaduras e fraturas de fêmur e coluna.
  \item \textbf{Lobotomia Transorbital (Procedimento de Freeman):} Intervenção cirúrgica introduzida nos anos 1940 em que uma haste metálica (semelhante a um picador de gelo) era inserida pelo canto interno do olho até o cérebro para seccionar conexões do lobo frontal.
  \item \textbf{Caquexia:} Estado de desnutrição crônica extrema e definhamento muscular profundo.
  \item \textbf{Linhas de Beau:} Sulcos transversais nas unhas causados pela interrupção temporária do crescimento ungueal durante episódios graves de fome ou infecção.
  \item \textbf{Icterícia Unilateral:} Coloração amarelada restrita a um único olho, achado pericial anômalo narrado no exame ficcional de Joaquim (Paciente 1.260).
\end{itemize}

\vspace{0.8em}
\noindent\rule{\textwidth}{0.5pt}
\vspace{0.6em}

\noindent\textbf{3. TERMOS PERICIAIS E ARCHIVÍSTICOS (2026)}

\vspace{0.4em}
\begin{itemize}[itemsep=0.2em]
  \item \textbf{Leucótomo:} Instrumento cirúrgico metálico pontiagudo usado em lobotomias.
  \item \textbf{Lâmina Cribriforme:} Placa óssea perfurada no teto da cavidade nasal/ocular que dá acesso à base do cérebro.
  \item \textbf{Pacto Ficcional Ergódico:} Dispositivo narrativo em que a estrutura física da obra e as decisões do leitor simulam o próprio contágio descrito no diário.
\end{itemize}

\endgroup
\clearpage
